\section{Background and relation to existing work}\label{sec:literature}
The literature approaches bubbles through several different observable
features: prices relative to cash flows, explosive growth, volatility that
increases with the price level, and run-ups followed by crashes. Each approach
starts from an economic mechanism and then asks which implications can be
measured. The present research connects the strict-local-martingale approach
to practical monitoring and to traded contracts. Its closest predecessors
therefore include both mathematical valuation papers and empirical detection
methods.

\subsection{From fundamental value to a pricing defect}
In the classical present-value formulation, an asset's fundamental value is
the discounted value of its future cash flows. A bubble is an additional
price component consistent with the pricing equation but not explained by
those cash flows. Empirical tests must consequently specify the cash-flow
process and the way investors discount it. G\"urkaynak's survey explains why
this creates an identification problem: evidence attributed to a bubble can
also be explained by changes in fundamentals or discounting
\cite{gurkaynak2005econometric}. A striking price path alone does not resolve
that comparison.

The strict-local-martingale literature provides a precise continuous-time
valuation framework. Cox and Hobson study prices that are local martingales
under a pricing measure but lose expectation, including local-volatility
and feedback models \cite{cox2005local}. Jarrow, Protter, and Shimbo place
such behavior in a theory of asset-price bubbles and market completeness
\cite{jarrow2007complete}. In the finite-horizon, discounted setting used
here, the object is the defect in \eqref{eq:defect}. This makes the
mathematical question explicit: under which models and trading rules does
current value exceed expected terminal value under the specified pricing law?

The specification of that law matters in an incomplete market. Herdegen
and Schweizer define fundamental value through superreplication and develop
a symmetric treatment of assets and numeraires. Their strong-bubble concept
differs from a bubble defined under one chosen equivalent local martingale
measure; their Example~5.5 explicitly exhibits different martingale behavior
under two such measures \cite{herdegen2016strong}. Our results concern the
fixed-measure defect. The contribution of the identification analysis is to
make the implications for observation precise: tail completions and changes
of pricing measure can preserve the available price information while
changing that defect.

\subsection{What volatility can reveal about a bubble}
Jarrow, Kchia, and Protter offer the closest methodological starting point
for the signals in this manuscript \cite{jarrow2011detect}. For a suitable
driftless diffusion $\dd S_t=\sigma(S_t)\dd W_t$, strictness is governed by
the upper-tail integral
\[
 \int_c^\infty \frac{x}{\sigma^2(x)}\dd x.
\]
The economic idea is that sufficiently rapid growth of volatility at high
prices can generate a pricing defect. Their detection procedure estimates
the diffusion coefficient from observed prices and then assesses its tail
behavior. They explicitly identify the extrapolation problem and propose
parametric and reproducing-kernel methods for extending the estimate beyond
the observed range.

Nonparametric diffusion estimation supplies the local measurement tools.
Florens-Zmirou studies estimation from discretely observed paths;
Bandi and Phillips develop nonparametric methods using occupation-time
information \cite{florenszmirou1993estimating,bandi2003fully}. Repeated visits
to a price region provide information about volatility there. More frequent
observations improve that local information, but need not reveal how the
coefficient behaves at prices never reached. Obayashi, Protter, and Yang
extend the detection approach toward dating bubble episodes: they construct
daily signals from intraday observations and smooth the noisy classifications
with a hidden Markov model \cite{obayashi2015lifetime}. Thus both
volatility-based detection and the use of intraday data to produce a daily
monitor have clear precedents.

Our theory specifies exactly what remains undetermined when the coefficient
is known only on a visited interval, including the lower endpoint that
matters for rebalanced wealth. The testing results then turn this observation
into a statistical statement: within the stated model classes, a
true-martingale completion can approximate the observed law of a strict
local martingale arbitrarily well. These are precise formulations of an
existing identification concern, with consequences for what assumptions a
useful test must impose.

The empirical methods take a more local target. They measure whether
\emph{return variance} rises with the price level after accounting for a
declared background scale. For intuition, if $\sigma(x)=a x^p$, then return
variance is $a^2x^{2(p-1)}$. The positive exponent estimated by a signal is
therefore connected to the classical tail mechanism, while remaining an
observable feature of the fitted window. Our methodological contribution is
the construction of that measurement under changing market volatility,
together with its calibration and comparative evidence. It is a different
use of the volatility insight from extrapolating a fitted coefficient to
infinity.

\subsection{Explosive prices, episode dating, and crash risk}
Another influential approach asks whether prices exhibit temporarily
explosive autoregressive dynamics. In a simple representation, the question
is whether the coefficient in $y_t=\rho y_{t-1}+u_t$ exceeds one over some
interval. The generalized supremum augmented Dickey--Fuller test scans
windows with varying starting and ending dates. This flexibility lets it
detect multiple episodes that a test tied to one starting date may miss;
the associated backward-window statistics support episode dating
\cite{phillips2012multiple,phillips2013realtime}. Panel methods can gain
information by grouping series with similar explosive behavior
\cite{liu2022panel}.

These tests give a natural practitioner benchmark because they formalize
the familiar pattern of accelerating prices. Their statistical target is
different from the tail condition above: explosive conditional-mean
dynamics and price-dependent conditional variance can appear separately.
This distinction makes the comparison informative. It lets us ask whether
the proposed variance signals detect patterns that an explosive-price
statistic tends to miss, and whether those patterns occur in different
historical episodes.

Volatility adjustment is already an important part of the explosive-testing
literature. Kurozumi, Skrobotov, and Tsarev construct a time transformation
to handle nonstationary volatility and recover a pivotal limiting null
distribution \cite{kurozumi2021time}. Our findings about the calibration
of GSADF belong in that context. The empirical comparisons here evaluate
the specified GSADF implementation and its calibrated thresholds; they
are not comparisons against every volatility-robust extension. The
distinctive evidence is the side-by-side behavior of a variance signal
and an explosive-price benchmark under common simulation and historical
designs.

A further strand studies what happens after observed run-ups. Greenwood,
Shleifer, and You find that sharp increases do not generally imply low
subsequent average returns, but do imply higher crash probabilities;
characteristics such as volatility, turnover, and issuance help distinguish
the outcomes \cite{greenwood2017bubbles}. This is the relevant precedent
for our supplementary drawdown comparison. The intraday result adds
information beyond the GSADF score in the reported forecasting design.
Its contribution is that particular incremental comparison, rather than
the general proposition that volatility or run-up characteristics can
predict crashes.

\subsection{Separating price dependence from background volatility}
Financial volatility changes over time and responds asymmetrically to
returns. Glosten, Jagannathan, and Runkle provide an influential model of
that asymmetry \cite{glosten1993relation}. Such dynamics matter here
because an observed price--variance association can arise through the
joint evolution of prices and a separate volatility state. A useful
measurement procedure must specify which variation its background model
explains and which variation remains available to the price-level term.

The connection to stochastic-volatility bubble theory is also substantive.
In Sin's model, positive alignment of price and volatility shocks can
produce strict-local-martingale behavior; the same model also illustrates
the role of the pricing measure \cite{sin1995complications}. Price--volatility
association is consequently of interest beyond one-dimensional local
volatility. The simulations in this research examine how the signals respond
to both structural level dependence and correlated volatility shocks.
That response helps explain what the statistic measures across model
families.

The daily construction uses a predictable market volatility scale and
compares positive and nonpositive price-level dependence with the same
nuisance specification. The intraday construction anchors the price
within each session and updates the background state between sessions.
These choices preserve a distinction between movements inside the
measurement window and changes in the background scale. Their value is
assessed through simulation, historical alarm patterns, and calibration
transfer. This is the central methodological addition to the earlier
volatility-estimation approach.

The statistical tools used to assess evidence also have established
antecedents. Split-likelihood universal inference supplies a way to
construct finite-sample evidence by estimating an alternative on separate
data and bounding the null likelihood \cite{wasserman2019universal}.
The e-BH procedure controls false discovery rate from valid e-values
under arbitrary dependence \cite{wang2020false}. Rank calibration uses
exchangeability to compare a new score with reference scores
\cite{tibshirani2019conformal}. We apply these ideas to specified
observation models and distinguish their guarantees from the empirical
reference calibration used by the practical scores. The findings about
how calibration changes sensitivity and historical alarms are part of
the contribution, alongside the score designs themselves.

\subsection{Investment, rebalancing, and leveraged funds}
A bubble in an asset raises a separate question about the wealth generated
by trading it. Herdegen and Herrmann connect strict local martingales with
a generalized Johansen--Ledoit--Sornette crash model and solve a power-utility
investment problem. Under their assumptions, investors with relative risk
aversion above one continue to hold the bubble when instantaneous expected
returns are positive \cite{herdegen2013strict}. This gives a useful
investment perspective: the existence of a bubble does not by itself
determine an optimal position. Their crash model also connects the
strict-local-martingale literature to the separate tradition of modeling
log-periodic behavior before crashes.

Our investment question concerns a prescribed rule: keep a constant fraction
of wealth in the asset. The tools are established martingale and diffusion
criteria. Mijatovi\'c and Urusov give the stochastic-exponential criterion
we apply, while Karatzas and Ruf analyze the auxiliary drift family that
arises in the calculation \cite{mijatovic2010martingale,karatzas2015distribution}.
The result developed here classifies the exposure rule across both price
endpoints. In particular, continuous buying into an asset approaching zero
can create strictness in wealth even when the asset is a true martingale.
The comparison with finite trading makes the economic importance of that
boundary mechanism concrete.

For leveraged exchange-traded funds, Avellaneda and Zhang establish the
path dependence of multi-period returns and their relationship to realized
variance, financing, and expenses \cite{avellaneda2009path}. That work
explains why a fund's cumulative return differs from a fixed multiple of
the underlying's cumulative return. Part~\ref{part:funds} addresses a
complementary issue: how finite rebalancing, nonnegativity, and liquidation
affect the martingale classification. The empirical tracking-cost evidence
then places those idealized wealth processes beside observed funds. The
established variance-drag explanation and the endpoint classification
answer different questions about the same trading rule.

\subsection{What options and constituents add}
The option-pricing consequences of strict local martingales are already
central to Cox and Hobson: familiar put--call parity and European--American
relations can change, and admissible trading strategies affect valuation
\cite{cox2005local}. Carr, Fisher, and Ruf develop a joint-superreplication
framework that produces a parity-restoring price \cite{carr2012hedging};
Bayraktar, Kardaras, and Xing study related strict-local-martingale and
American-option phenomena \cite{bayraktar2011strict}. These papers supply
the valuation background for Part~\ref{part:options}.

Our contribution in that part is to organize the information supplied by
actual contract observations. A quote becomes informative about the defect
only after its pricing convention and exercise rights have been specified.
The resulting inequalities distinguish an expectation-priced European
option, a parity-restored value, and an American exercise value. This
provides a practical way to read option evidence and to understand why
different contract conventions can reveal or absorb the same defect.

Constituent information offers another route beyond an aggregate price.
The underlying change-of-measure and explosion ideas come from the
strict-local-martingale literature, including Kardaras, Kreher, and
Nikeghbali \cite{kardaras2015strict}. We use them to express constituent
defects through weights stopped at the basket's lifetime and their
conditional expectations. The useful addition is an explicit attribution
formula and an account of the extra information its recovery requires.
Together, the fund, option, and constituent results extend the question
from detecting a feature of one price series to understanding the
exposures and contracts an investor actually holds.
